\documentclass[11pt]{article}
\usepackage[margin=1in]{geometry}
\usepackage[T1]{fontenc}
\usepackage{lmodern,microtype,hyperref,enumitem,amsmath,amssymb,xcolor,booktabs}
\usepackage[numbers,sort&compress]{natbib}
\hypersetup{colorlinks=true,urlcolor=blue,citecolor=blue,linkcolor=blue}
\title{Q10: Comprehension Check}
\author{DLS 2026 --- Kukla}
\date{October 2026}
\begin{document}
\maketitle
\textbf{Comprehension check.} A strong model correctly derives a calculus result but says the word
``banana'' has two a's. Give the most likely technical reason for the counting error, and name a change to
the \emph{input or harness} (not the model weights) that would fix it.

\vspace{0.8cm}\hrule\vspace{4pt}
{\small\textbf{Answer.} The model does not see individual letters---it processes sub-word \emph{tokens}, so
``banana'' is one or a few tokens with no explicit per-letter inventory, and counting a specific letter
requires structure the representation does not expose. Fixes that leave the weights untouched: reframe the
input so the letters are individually visible (e.g.\ ask it to first spell the word out letter-by-letter,
``b-a-n-a-n-a'', then count), or give it a tool (run a tiny code snippet / character counter). Either
exposes the sub-token structure the raw tokenization hides. (``banana'' in fact has three a's---exactly
the kind of slip the reframing prevents.)}
\end{document}
